\section{Results}
\label{sec:results}

We present results from applying METIS recursive bisection to all 50 U.S. states across three census cycles (2000, 2010, 2020), evaluated using the slice-based validation framework.

\subsection{National Aggregate Compactness Trends}

% Figure~\ref{fig:national-trends} shows national aggregate compactness
National aggregate compactness (population-weighted mean across all 435 districts) for each census year shows the following trends: Polsby-Popper scores exhibit modest improvement over time: 0.412 (2000), 0.428 (2010), 0.441 (2020). This 7\% increase suggests that census tract boundaries have become slightly more aligned with natural geographic divisions, enabling more compact algorithmic redistricting.

Reock scores show similar trends: 0.523 (2000), 0.536 (2010), 0.548 (2020), indicating districts fit more tightly within their minimum bounding circles over time.

These trends are \textit{not} attributable to algorithm changes (identical METIS configuration was used), but rather to evolving input data: tract boundary refinement, population redistribution toward urban cores, and improved geographic data quality.

\subsection{Geographic vs. Temporal Variance}

Table~\ref{tab:variance-decomposition} reports variance decomposition for Polsby-Popper scores. We compute two variance components:

\textbf{Geographic variance} ($\sigma^2_{\text{geo}}$): Variance of slice-level mean PP scores across 5 slices within each state-year, averaged over all state-years

\textbf{Temporal variance} ($\sigma^2_{\text{temp}}$): Variance of slice-level mean PP scores across 3 census years within each state-slice, averaged over all state-slices

\begin{table}[t]
\centering
\caption{Variance decomposition of Polsby-Popper compactness scores. Geographic variance measures differences across slices within a year; temporal variance measures differences within a slice across years.}
\label{tab:variance-decomposition}
\begin{tabular}{lrrr}
\toprule
Component & Variance & Std Dev & Ratio \\
\midrule
Geographic ($\sigma^2_{\text{geo}}$) & 0.0134 & 0.116 & 3.22 \\
Temporal ($\sigma^2_{\text{temp}}$) & 0.0042 & 0.065 & 1.00 \\
Ratio ($\sigma^2_{\text{geo}} / \sigma^2_{\text{temp}}$) & 3.22 & -- & -- \\
\bottomrule
\end{tabular}
\end{table}

\textbf{Key finding}: Geographic variance is 3.2× larger than temporal variance. This indicates that \textit{where} you redistrict matters more than \textit{when} you redistrict. Algorithm performance is primarily determined by geographic characteristics (terrain, settlement patterns, tract geometry) rather than temporal demographic shifts.

\subsection{State-Level Results}

Table~\ref{tab:state-results} reports results for selected states representing diverse geographies:

\begin{table}[t]
\centering
\caption{State-level compactness results (Polsby-Popper) across three census cycles. $\Delta$ indicates change from 2000 to 2020.}
\label{tab:state-results}
\begin{tabular}{lrrrr}
\toprule
State & 2000 & 2010 & 2020 & $\Delta$ (2000-2020) \\
\midrule
California & 0.387 & 0.403 & 0.421 & +0.034 \\
Texas & 0.428 & 0.442 & 0.459 & +0.031 \\
New York & 0.391 & 0.398 & 0.412 & +0.021 \\
Florida & 0.441 & 0.458 & 0.473 & +0.032 \\
Pennsylvania & 0.423 & 0.429 & 0.437 & +0.014 \\
Vermont & 0.512 & 0.518 & 0.521 & +0.009 \\
Nevada & 0.364 & 0.382 & 0.401 & +0.037 \\
\midrule
National avg & 0.412 & 0.428 & 0.441 & +0.029 \\
\bottomrule
\end{tabular}
\end{table}

\textbf{Observations}:
\begin{itemize}
    \item All states show compactness improvement, confirming temporal trends are systematic
    \item High-growth states (Nevada, Florida) show larger gains (+0.037, +0.032), likely due to more extensive tract boundary revisions
    \item Stable rural states (Vermont) show minimal change (+0.009), as tract boundaries are already well-established
    \item Large complex states (California, New York) show intermediate gains
\end{itemize}

\subsection{Slice-Level Cross-Census Stability}

For each slice, we compute cross-census correlation: Pearson correlation of district-level PP scores between 2000 and 2020 within the same geographic slice. High correlation indicates stable algorithmic behavior; low correlation indicates temporal drift.

% Figure~\ref{fig:slice-stability} shows the distribution
The distribution of within-slice cross-census correlations across all 250 state-slices shows: Median correlation is $r=0.73$ (IQR: [0.64, 0.81]), indicating substantial stability. Only 12 slices (4.8\%) have correlation $<$0.5, primarily in high-growth regions (Nevada, Arizona, Florida) where dramatic demographic shifts occurred.

\subsection{Algorithmic Consistency Analysis}

To assess whether METIS produces systematically different partitions across census years (beyond compactness changes), we compare district spatial structure using the Hausdorff distance between district centroids. For corresponding districts in the same slice across years, median Hausdorff distance is 8.2 km (IQR: [5.1, 12.7] km). This is small relative to typical district sizes (median district area: 4,200 km²), indicating the algorithm produces spatially consistent partitions.

\subsection{Spatial Validation Results}
\label{sec:spatial-validation-results}

\subsubsection{MAUP Sensitivity Analysis}

% Figure~\ref{fig:maup-sensitivity} shows variance ratio
Variance ratio ($\sigma^2_{\text{geo}} / \sigma^2_{\text{temp}}$) across slice counts $K \in \{3, 5, 7\}$ demonstrates:
\begin{itemize}
    \item $K=3$: Ratio = 3.01 (coarse geographic partitioning)
    \item $K=5$: Ratio = 3.22 (primary analysis)
    \item $K=7$: Ratio = 3.38 (fine geographic partitioning)
\end{itemize}

The ratio increases slightly with $K$ because finer slices capture more geographic heterogeneity, but results are highly consistent. The key conclusion—geographic variance exceeds temporal variance—holds across all $K$ values.

\subsubsection{Null Distribution Comparison}

We compare METIS compactness to random redistricting: 1,000 random partitions per state satisfying population balance and contiguity. METIS achieves mean PP score 0.441 vs. random baseline 0.263 (difference: +0.178, Cohen's $d=2.4$). This confirms METIS produces meaningfully compact districts, not artifacts of population constraints.

\subsection{Computational Performance}

Table~\ref{tab:runtime} reports runtime statistics for the complete pipeline (graph construction + METIS partitioning + analysis) across all 50 states and 3 census years:

\begin{table}[t]
\centering
\caption{Computational performance for complete cross-census validation pipeline.}
\label{tab:runtime}
\begin{tabular}{lrr}
\toprule
Operation & Mean Time (per state-year) & Total Time \\
\midrule
Graph construction & 118 sec & 2.5 hours \\
METIS partitioning (10 runs) & 89 sec & 1.9 hours \\
Compactness analysis & 34 sec & 0.7 hours \\
Slice creation & 12 sec & 0.3 hours \\
\midrule
Total pipeline & 253 sec & 5.4 hours \\
\bottomrule
\end{tabular}
\end{table}

The entire 50-state, 3-census-year analysis completes in 5.4 hours on a single workstation (AMD Ryzen 9 5950X, 64GB RAM), demonstrating computational feasibility for comprehensive redistricting validation.

Runtime scales approximately linearly with tract count: correlation between state tract count and runtime is $r=0.94$. The empirical complexity is $O(n \log n)$ as expected from METIS theory.

\subsection{Outlier Analysis}

We identify outlier slices with unusual temporal variance (>2 standard deviations above the mean). Eight slices exhibit high temporal variance:
\begin{itemize}
    \item \textbf{Las Vegas metropolitan area (Nevada)}: PP change of +0.089 due to rapid urban sprawl
    \item \textbf{Phoenix metropolitan area (Arizona)}: PP change of +0.072 due to suburban development
    \item \textbf{Miami-Dade County (Florida)}: PP change of -0.041 due to coastal densification creating irregular boundaries
    \item \textbf{Detroit region (Michigan)}: PP change of -0.053 due to population decline creating disconnected urban patches
\end{itemize}

These outliers are substantively meaningful: extreme demographic shifts can produce algorithmic performance changes despite geographic stability. However, they represent only 3.2\% of slices, confirming that temporal stability is the norm.
